% భాగం: సమితులు, సంబంధాలు, ప్రమేయాలు
% అధ్యాయం: సమితుల పరిమాణం

\documentclass[../../../include/open-logic-chapter]{subfiles}

\begin{document}

\olchapter{sfr}{siz}{సమితుల పరిమాణం}

\begin{editorial}
మూలకాలను జాబితా చేయడం, లెక్కించదగిన సమితులు, లెక్కించలేని సమితులు ఈ అధ్యాయంలోని అంశాలు. కొన్ని విభాగాలకు రెండు రూపాలున్నాయి. ఒకటి ప్రాథమిక రూపం. సమితి సిద్ధాంతాన్ని మరింత వివరంగా అధ్యయనం చేసేటప్పుడు మరో ఉన్నత స్థాయి రూపాన్ని వాడవచ్చు. ఈ దస్త్రంలో ప్రాథమిక రూపాలు మాత్రమే ఉన్నాయి. ఉన్నత స్థాయి రూపాలు \verb|size-of-sets-complete|లో ఉన్నాయి.
\end{editorial}

\olimport{introduction}

\olimport{enumerability}

\olimport{zig-zag}

\olimport{pairing}

\olimport{pairing-alt}

\olimport{non-enumerability}

\olimport{reduction}

\olimport{equinumerous-sets}

\olimport{comparing-size}

\olimport{schroder-bernstein}

\OLEndChapterHook

\end{document}
